% ---------------------------------------------------------------------
%  ABSTRACT — tradução do resumo; manter os dois sincronizados.
% ---------------------------------------------------------------------
\begin{resumo}[Abstract]
\begin{otherlanguage*}{english}
Quadrotors are underactuated, nonlinear, coupled and open-loop unstable
plants, and their flight depends on an attitude control loop running in real
time. In open-source autopilots this loop is built from cascaded
proportional-integral-derivative (PID) controllers, tuned empirically or
automatically rather than from a model of the vehicle. This work designs
three model-based attitude controllers --- LQI, LQR with an extended state
observer and LQG ---, implements them inside the ArduPilot firmware through
the \codigo{AC\_CustomControl} framework and compares them with the native
PID controller, tuned by AutoTune. The comparison runs in software-in-the-loop
(SITL) simulation, in which the compiled firmware flies against a flight
dynamics model. The vehicle model is derived from the Newton-Euler equations,
linearised around hover and checked against the simulator, in the time and
frequency domains. The LQI and the LQR with observer use the native
bandwidth; the LQG uses a Kalman filter that models the firmware's
measurement path and is designed to pass no more gyroscope noise to the
motors than the native controller. The tests apply to all controllers, under
the same conditions, attitude steps and torque disturbances shaped as
impulses, short steps, sustained steps, frequency sweeps and stochastic
processes. There are also scenarios with gyroscope noise, model error, motor
failure, saturation and large angles, in a batch of 558 flights with frozen
firmware. In roll and pitch, the LQI and the LQR with observer do not
outperform the native controller in peak disturbance rejection, and the LQI
is equivalent to the native cascade with different gains. The LQG is the
only one that reduces the peak, by 7\,\% to 11\,\%, also with gyroscope noise
and in repeated flights. An ablation shows that this advantage comes from
modelling the gyroscope filter and the error of the estimated angle. In yaw,
all proposed controllers remove the native controller's slow recovery, from
22~s to 0.7 to 2.3~s, but a retuned PID does so as well; the peak only drops
with more bandwidth and more noise in the control effort or, with the LQG,
by 7\,\% to 8\,\%. The advantage in attitude step tracking comes from
feedforward by model inversion, which the PID itself could use. The limits
found stem more from the available information --- the gyroscope filter and
noise and the low yaw authority --- than from the control law.

\vspace{\onelineskip}

\noindent
\textbf{Keywords}: quadrotor; attitude control; ArduPilot; linear quadratic
regulator; Kalman filter; software-in-the-loop simulation.
\end{otherlanguage*}
\end{resumo}
